% Source: https://artofproblemsolving.com/community/c3h2940557_the_movement_of_thex10andx20when_a_parabola_goes_up_and_down
% Stable identity: aops-topic-2940557-q01
The intercepts are $x=\dfrac{-b\pm\sqrt{\Delta}}{2a}$ with $\Delta=b^2-4a(c+k)>0$. Their midpoint is $-\dfrac b{2a}$, the axis of symmetry, which does not depend on $k$.
\begin{asy}
usepackage("amsmath");
size(4cm);
defaultpen(fontsize(10pt));
import graph;real f(real x){return x^2-2*x-2;}real g(real x){return x^2-2*x;}
draw(graph(f,-1.6,3.6));draw(graph(g,-1.2,3.2),dashed);draw((-2,0)--(4,0),Arrow);draw((1,-3.5)--(1,4),dotted);label("$x=-\dfrac b{2a}$",(1,4),N);
\end{asy}

The distance between them is
\[\boxed{\dfrac{\sqrt{b^2-4a(c+k)}}{|a|}}.\]
So the intercepts move symmetrically, but not at a constant rate: the distance is a square-root function of $k$.
